\subsection{2.10 Die physikalische Interpretation der Projektion}

Mit den numerischen Untersuchungen der vorangegangenen Kapitel wurde ein Punkt erreicht, an dem sich die mathematischen Ergebnisse erstmals mit den ursprünglichen philosophischen Grundannahmen des Ontophänomenalismus (OP) verbinden ließen.

Bis hierher hatten die Simulationen im Wesentlichen Folgendes gezeigt:

\begin{enumerate}
\item Der rohe OPP besitzt keine intrinsische vierdimensionale Struktur.
\item Verschiedene Projektionsverfahren erzeugen unterschiedliche effektive Dimensionen.
\item Nur wenige Projektionen führen stabil in den Bereich
\[
d_{\mathrm{eff}}\approx4.
\]
\item Diffusion Maps und Kernel PCA erwiesen sich als die robustesten Kandidaten.
\end{enumerate}

Damit stellte sich nun die entscheidende Frage:

\[
\boxed{
\text{Warum sollte gerade eine Projektion physikalische Realität erzeugen?}
}
\]

%%%%%%%%%%%%%%%%%%%%%%%%%%%%%%%%%%%%%%%%%%%%%%%%%%%%%%%%%%%%%%

\section*{Von der Mathematik zurück zur Physik}

Bis zu diesem Zeitpunkt hätte man sämtliche Resultate auch rein mathematisch interpretieren können.

Ein Mathematiker würde sagen:

\begin{quote}
„Es existieren verschiedene Möglichkeiten, einen hochdimensionalen Datensatz niedrigdimensional darzustellen.“
\end{quote}

Die PST geht jedoch einen Schritt weiter.

Sie behauptet:

\[
\boxed{
\text{Die beobachtbare Physik ist genau eine solche Projektion.}
}
\]

Damit wird aus einem mathematischen Werkzeug plötzlich ein physikalischer Mechanismus.

%%%%%%%%%%%%%%%%%%%%%%%%%%%%%%%%%%%%%%%%%%%%%%%%%%%%%%%%%%%%%%

\section*{Die Analogie zur Sinneswahrnehmung}

Bereits während der Diskussion entstand eine wichtige Analogie.

Der Mensch erlebt niemals die vollständige physikalische Realität.

Unsere Wahrnehmung ist grundsätzlich begrenzt.

Beispiele:

\begin{itemize}

\item

sichtbares Licht umfasst nur einen winzigen Bereich des elektromagnetischen Spektrums,

\item

Schall wird nur innerhalb eines kleinen Frequenzfensters wahrgenommen,

\item

viele Tiere besitzen vollkommen andere Wahrnehmungsräume.

\end{itemize}

Keine dieser Beschränkungen verändert die Realität selbst.

Sie verändern ausschließlich die beobachtete Realität.

Formal:

\[
\boxed{
\text{Beobachtung}
=
\text{Projektion der Realität}.
}
\]

%%%%%%%%%%%%%%%%%%%%%%%%%%%%%%%%%%%%%%%%%%%%%%%%%%%%%%%%%%%%%%

\section*{Übertragung auf den OPP}

Die gleiche Idee lässt sich unmittelbar auf den OPP übertragen.

Der OPP besitzt möglicherweise sehr viele Freiheitsgrade.

Ein realer Beobachter verfügt jedoch nur über endliche Informationskapazität.

Daher kann niemals der vollständige OPP beobachtet werden.

Es erscheint lediglich

\[
\mathcal{P}(\mathrm{OPP}).
\]

Die Raumzeit wäre somit nicht Bestandteil des OPP,

sondern Bestandteil der Beobachtung.

%%%%%%%%%%%%%%%%%%%%%%%%%%%%%%%%%%%%%%%%%%%%%%%%%%%%%%%%%%%%%%

\section*{Die Emergenz der Raumzeit}

Hier ergibt sich ein völlig neues Bild.

Nicht

\[
\text{Raum}
\]

ist fundamental,

sondern

\[
\text{Relation}.
\]

Nicht

\[
\text{Zeit}
\]

ist fundamental,

sondern

\[
\text{Informationsfluss}.
\]

Nicht

\[
\text{Materie}
\]

ist fundamental,

sondern

\[
\text{stabile Attraktoren}.
\]

Die klassische Physik erscheint lediglich als eine spezielle Beschreibungsebene.

%%%%%%%%%%%%%%%%%%%%%%%%%%%%%%%%%%%%%%%%%%%%%%%%%%%%%%%%%%%%%%

\section*{Lorentz-Invarianz}

Während der Diskussion entstand ein weiterer wichtiger Gedanke.

Falls Raumzeit tatsächlich projektional entsteht,

dann könnte auch die Lorentz-Invarianz eine Konsequenz derselben Projektion sein.

Bisher wurde meist angenommen,

dass Lorentz-Symmetrie fundamental sei.

Die PST schlägt stattdessen vor:

\[
\boxed{
\text{Lorentz-Invarianz ist eine Eigenschaft der Projektion.}
}
\]

Nicht der OPP muss Lorentz-invariant sein,

sondern ausschließlich die beobachtbare Raumzeit.

%%%%%%%%%%%%%%%%%%%%%%%%%%%%%%%%%%%%%%%%%%%%%%%%%%%%%%%%%%%%%%

\section*{Quantisierung}

Ein ähnlicher Gedanke betrifft die Quantenphysik.

Bereits früher wurde diskutiert,

dass Quantisierung möglicherweise nicht fundamental ist.

Falls jeder Beobachter nur eine endliche Informationsauflösung besitzt,

dann entstehen zwangsläufig diskrete Messwerte.

Formal:

\[
\boxed{
\text{Quantisierung}
=
\text{Projektionsartefakt}.
}
\]

Dies würde erklären,

warum kontinuierliche mathematische Strukturen diskrete Messergebnisse liefern.

%%%%%%%%%%%%%%%%%%%%%%%%%%%%%%%%%%%%%%%%%%%%%%%%%%%%%%%%%%%%%%

\section*{Der Beobachter wird physikalisch}

In vielen physikalischen Theorien erscheint der Beobachter erst am Ende.

Die PST kehrt diese Reihenfolge um.

Der Beobachter wird selbst Bestandteil der fundamentalen Struktur.

Nicht deshalb,

weil Bewusstsein die Realität erschafft,

sondern weil jede Beobachtung zwangsläufig eine Projektion darstellt.

Die physikalische Welt ist daher immer eine beobachtbare Welt.

%%%%%%%%%%%%%%%%%%%%%%%%%%%%%%%%%%%%%%%%%%%%%%%%%%%%%%%%%%%%%%

\section*{Die Verbindung zum OP}

An dieser Stelle schließen sich erstmals Ontophänomenalismus und PST.

Der OP behauptet:

\[
\text{Existenz ist fundamental.}
\]

Die PST ergänzt:

\[
\text{Bestimmtheit entsteht durch Projektion.}
\]

Damit ergibt sich:

\[
\boxed{
\text{Existenz}
\rightarrow
\text{OPP}
\rightarrow
\text{Projektion}
\rightarrow
\text{Bestimmtheit}
\rightarrow
\text{Physik}.
}
\]

Genau diese Struktur war bereits in den philosophischen Diskussionen über Existenz, Differenz und Bestimmtheit vorbereitet worden.

%%%%%%%%%%%%%%%%%%%%%%%%%%%%%%%%%%%%%%%%%%%%%%%%%%%%%%%%%%%%%%

\section*{Der neue Forschungsansatz}

Aus dieser Interpretation ergibt sich eine völlig neue Forschungsstrategie.

Nicht mehr die Frage

\[
\text{„Wie entsteht Raumzeit?“}
\]

steht im Mittelpunkt.

Sondern:

\[
\boxed{
\text{Welche Projektionsoperatoren besitzen dieselben Eigenschaften wie reale Beobachter?}
}
\]

Damit wird die Suche nach geeigneten Projektionsverfahren zu einer physikalischen Fragestellung.

%%%%%%%%%%%%%%%%%%%%%%%%%%%%%%%%%%%%%%%%%%%%%%%%%%%%%%%%%%%%%%

\section*{Eine überraschende Konsequenz}

Während der numerischen Arbeiten entstand schließlich noch eine weitere Vermutung.

Vielleicht existiert gar keine einzige ausgezeichnete Projektion.

Möglicherweise existiert vielmehr eine ganze Klasse physikalisch zulässiger Projektionen.

Unterschiedliche Beobachter würden dann unterschiedliche, aber kompatible Raumzeiten erzeugen.

Dies erinnert unmittelbar an die Spezielle Relativitätstheorie.

Auch dort existiert kein ausgezeichnetes Inertialsystem.

%%%%%%%%%%%%%%%%%%%%%%%%%%%%%%%%%%%%%%%%%%%%%%%%%%%%%%%%%%%%%%

\section*{Offene Fragen}

Trotz der bisherigen Erfolge bleiben zahlreiche fundamentale Fragen offen.

Insbesondere:

\begin{itemize}

\item

Warum bevorzugen Diffusion Maps und Kernel PCA ähnliche Dimensionen?

\item

Welche mathematischen Eigenschaften muss ein physikalischer Projektionsoperator besitzen?

\item

Wie entstehen Eichsymmetrien innerhalb einer Projektion?

\item

Kann die Lorentz-Metrik vollständig aus einem Projektionsprinzip folgen?

\item

Welche Rolle spielt der Beobachter mathematisch genau?

\end{itemize}

Diese Fragen bilden den Übergang vom numerischen Teil der Arbeit zur eigentlichen physikalischen Theorieentwicklung.

%%%%%%%%%%%%%%%%%%%%%%%%%%%%%%%%%%%%%%%%%%%%%%%%%%%%%%%%%%%%%%

\section*{Zusammenfassung}

Die bisherigen Simulationen führten zu einem grundlegenden Perspektivwechsel.

Nicht der OPP besitzt notwendigerweise eine vierdimensionale Raumzeit.

Vielmehr entsteht Raumzeit als stabile Projektion eines wesentlich reichhaltigeren ontophänomenalen Phasenraums.

Damit verbindet die PST erstmals:

\begin{itemize}

\item die philosophischen Überlegungen des Ontophänomenalismus,

\item die numerischen Ergebnisse der OPP-Simulationen,

\item moderne Methoden der Dimensionsreduktion,

\item sowie die Idee einer beobachterabhängigen Emergenz der Physik.

\end{itemize}

Dieser Gedanke bildet die Grundlage aller folgenden Entwicklungen der PST und stellt den eigentlichen konzeptionellen Kern der Theorie dar.

